\chapter{Metody nieparametryczne i odporne: rozkład, rangi i obserwacje odstające}
\label{ch:nonparametric}
\paragraph{Powiązania i odsyłacze.} Dystrybuantę empiryczną omówiliśmy w rozdziale~\ref{ch:statystyki}, a bootstrap i testy permutacyjne w rozdziale~\ref{ch:resampling}.


\section{Co oznacza brak założenia parametrycznego?}
W modelu parametrycznym ustalamy rodzinę rozkładów
$\{P_\theta:\theta\in\Theta\}$, a w modelu nieparametrycznym
przestrzeń dopuszczalnych rozkładów jest znacznie większa,
np. obejmuje wszystkie rozkłady o skończonej wariancji.
Nie znaczy to, że metoda nieparametryczna nie wymaga
\emph{żadnych} założeń: test rangowy może wymagać wymienności,
estymacja gęstości -- gładkości, a test dla par -- odpowiednio
symetrycznego rozkładu różnic. Często wybieramy pomiędzy
\emph{różnymi hipotezami} i różnymi miarami efektu,
a nie tylko pomiędzy dwoma wzorami na wartość p.

\section{Dystrybuanta empiryczna i dokładność jej aproksymacji}
\begin{definicja}[Dystrybuanta empiryczna]
Niech $X_1,\ldots,X_n$ będą niezależne o jednakowej
dystrybuancie $F$. Definiujemy
$F_n(t)=n^{-1}\sum_{i=1}^n\ind{\{X_i\le t\}}$,
$t\in\R$. Dla ustalonego $t$ zmienna $nF_n(t)$
ma rozkład $\Bin(n,F(t))$.
\end{definicja}

\begin{twierdzenie}[Jednopunktowa kontrola dystrybuanty]
\label{thm:empirical-hoeffding}
Dla dowolnego $n\ge1$, $t\in\R$ i $\varepsilon>0$,
\[
 \Prob\bigl(|F_n(t)-F(t)|\ge\varepsilon\bigr)
 \le2\exp(-2n\varepsilon^2).
\]
\end{twierdzenie}
\begin{proof}
Zmienne $\ind{\{X_i\le t\}}$ są niezależne, przyjmują
wartości w $[0,1]$ i mają wspólną średnią $F(t)$.
Nierówność Hoeffdinga dla średniej z ograniczonych
zmiennych daje tezę.
\end{proof}

\begin{uwaga}[Co znaczy jednoczesny pas ufności?]
Nierówność powyższa dotyczy \emph{ustalonego} $t$.
Nie wolno jej bez uzasadnienia stosować równocześnie
do wszystkich $t\in\R$. Silniejsza nierówność
Dvoretzky'ego--Kiefera--Wolfowitza--Massarta głosi,
że dla próby niezależnej o jednakowym rozkładzie
\[
 \Prob\left(\sup_{t\in\R}|F_n(t)-F(t)|>\varepsilon\right)
 \le2\exp(-2n\varepsilon^2).
\]
Jest to głębsze twierdzenie, którego pełnego dowodu
tu nie przeprowadzamy. Daje jednoczesny pas
$F_n(t)\pm\sqrt{\log(2/\alpha)/(2n)}$, obcięty do
$[0,1]$, o pokryciu co najmniej $1-\alpha$.
Nie należy mylić go z punktowym przedziałem dla $F(t)$.
\end{uwaga}

\section{Mediana, kwantyle i odporność na skażenia}
\begin{definicja}[Funkcjonał kwantyla i punkt skażenia]
Dla dystrybuanty $F$ i $p\in(0,1)$ definiujemy
$Q_p(F)=\inf\{x\in\R:F(x)\ge p\}$.
Dla $0\le\varepsilon\le1$ i $z\in\R$ rozważamy
\emph{model skażenia}
$F_{\varepsilon,z}=(1-\varepsilon)F+\varepsilon\delta_z$,
gdzie $\delta_z$ oznacza rozkład punktowy w $z$.
\emph{Funkcją wpływu} funkcjonału $T$ w punkcie $z$ nazywamy,
jeśli granica istnieje,
\[
 \mathrm{IF}(z;T,F)=
 \lim_{\varepsilon\downarrow0}
 \frac{T(F_{\varepsilon,z})-T(F)}{\varepsilon}.
\]
Jest to miara lokalnej wrażliwości na niewielką domieszkę
obserwacji w punkcie $z$, nie zaś dokładny opis
arbitralnie dużego skażenia.
\end{definicja}

\begin{twierdzenie}[Funkcje wpływu średniej i kwantyla]
\label{thm:influence-mean-quantile}
Niech $F$ ma skończoną średnią $\mu$.
Dla funkcjonału $T(F)=\int_\R x\,\dd F(x)$
zachodzi $\mathrm{IF}(z;T,F)=z-\mu$.
Niech ponadto $q=Q_p(F)$ będzie jedynym kwantylem
rzędu $p$, a $F$ ma w otoczeniu $q$ ciągłą,
dodatnią gęstość $f$, tak że $F(q)=p$.
Wówczas dla każdego $z\ne q$
\[
 \mathrm{IF}(z;Q_p,F)=
 \frac{p-\ind{\{z\le q\}}}{f(q)}.
\]
Funkcja wpływu kwantyla jest zatem ograniczona
dla ustalonego $F$ i $p$, podczas gdy funkcja wpływu
średniej jest nieograniczona.
\end{twierdzenie}
\begin{proof}
Średnia przy skażeniu wynosi
$(1-\varepsilon)\mu+\varepsilon z$, co daje pierwszy wzór.
Dla $z\ne q$ przy małym $\varepsilon$ kwantyl
$q_\varepsilon=Q_p(F_{\varepsilon,z})$ pozostaje
w otoczeniu $q$ nie zawierającym $z$.
Rozwiązujemy zatem równanie
$(1-\varepsilon)F(q_\varepsilon)+
\varepsilon\ind{\{z\le q\}}=p$.
Różniczkując jednostronnie w $\varepsilon=0$,
przy $f(q)>0$, otrzymujemy
$-p+f(q)q'_0+\ind{\{z\le q\}}=0$.
To daje drugi wzór i wskazane ograniczenia.
\end{proof}

\begin{przyklad}[Jedna obserwacja odstająca i dwie różne odpowiedzi]
Rozważmy pomiary $9,10,10,11,12$. Średnia wynosi $10{,}4$,
a mediana $10$. Jeśli pomiar $12$ zostanie zastąpiony
przez $120$, nowa średnia wyniesie $32$, ale mediana
pozostanie równa $10$. Mediana ma wysoką odporność
na kilka skrajnych wartości, lecz przy regularnym modelu
normalnym może być mniej efektywna od średniej.
Wybór zależy od celu: jeśli mierzymy oczekiwany
\emph{łączny koszt}, średnia może być właściwym funkcjonałem
nawet dla rozkładu z ogonem; jeśli szukamy typowego
poziomu, mediana może być łatwiejsza do interpretacji.
\end{przyklad}

\begin{center}
\begin{tikzpicture}[x=.072cm,y=1.0cm,>=Latex]
\draw[->] (0,0)--(130,0) node[right] {pomiar};
\foreach \x in {0,20,40,60,80,100,120} {
 \draw (\x,-.06)--(\x,.06) node[below=4pt,font=\scriptsize] {\x};
}
\foreach \x in {9,10,11} {
 \fill[blue!70!black] (\x,0.40) circle[radius=2.7pt];
}
\fill[blue!70!black] (10,.62) circle[radius=2.7pt];
\fill[red!75!black] (120,.40) circle[radius=3.2pt];
\draw[red!70!black,very thick] (32,.18)--(32,.86)
 node[above,font=\scriptsize] {średnia $32$};
\draw[green!50!black,very thick] (10,.86)--(10,1.35)
 node[above,font=\scriptsize] {mediana $10$};
\node[anchor=west,font=\scriptsize] at (52,1.20)
 {Punkt czerwony zastępuje $12$ przez $120$};
\end{tikzpicture}

{\small\emph{Rysunek 26.1.} Jeden skrajny pomiar może silnie przesunąć średnią,\newline
nie zmieniając mediany. Punkty w pobliżu zera nachodzą na siebie w tej skali.}
\end{center}

\section{Testy rangowe dla dwóch niezależnych grup}
\begin{definicja}[Statystyka Manna--Whitneya]
\label{def:mann-whitney}
Dla niezależnych prób $X_1,\ldots,X_m$ i
$Y_1,\ldots,Y_k$ definiujemy
\[
 U=\sum_{i=1}^m\sum_{j=1}^k
 \left(\ind{\{X_i<Y_j\}}+
 \frac12\ind{\{X_i=Y_j\}}\right).
\]
Jeśli grupy są niezależne i każda jest próbą
z odpowiedniego rozkładu $F$ lub $G$, to
$\E U/(mk)=\Prob(X<Y)+\tfrac12\Prob(X=Y)$,
gdzie $X\sim F$, $Y\sim G$ są niezależne.
\end{definicja}

\begin{twierdzenie}[Rozkład rang i momenty $U$ pod hipotezą zerową]
\label{thm:mann-whitney-moments}
Niech wszystkie $m+k$ obserwacji będą niezależne
i mają wspólny \emph{ciągły} rozkład, przy $m,k\ge1$.
Wtedy
\[
 \E U=\frac{mk}{2},\qquad
 \Var(U)=\frac{mk(m+k+1)}{12}.
\]
Ponadto rozkład $U$ zależy tylko od $(m,k)$,
a nie od wspólnego ciągłego rozkładu obserwacji.
\end{twierdzenie}
\begin{proof}
Przez ciągłość wszystkie obserwacje są różne prawie na pewno.
W wyniku wymienności wspólnych obserwacji każda z
$\binom{m+k}{m}$ konfiguracji rang pierwszej grupy
jest jednakowo prawdopodobna. Jeśli $R_1,\ldots,R_m$
oznaczają jej rangi (od najmniejszej), to
$U=mk+m(m+1)/2-\sum_{i=1}^mR_i$ przy naszej
konwencji zliczania $X_i<Y_j$.
Z losowania bez zwracania $m$ rang z $1,\ldots,m+k$
wynika $\E\sum_iR_i=m(m+k+1)/2$.
Wariancja sumy rang jest równa
\[
 \frac{mk}{m+k-1}\cdot
 \frac{(m+k)^2-1}{12}
 =\frac{mk(m+k+1)}{12}.
\]
Odjęcie sumy rang od stałej nie zmienia wariancji.
Rozkład rang jest jednostajny, więc nie zależy od
wspólnej dystrybuanty.
\end{proof}

\begin{uwaga}[Mann--Whitney nie jest automatycznie testem median]
Hipoteza o identycznych rozkładach $F=G$ daje
\emph{dokładną} rangową kalibrację testu.
Przy rodzinie przesunięć o tym samym kształcie można
go interpretować jako test położenia.
Bez tego założenia różnica median nie jest równoważna
zmianie $\Prob(X<Y)$; grupy mogą mieć tę samą medianę,
ale inne wariancje i istotnie różne wartości statystyki $U$.
Gdy celem jest test \emph{równości średnich}, zamiast
Manna--Whitneya warto rozważyć test Welcha,
studentyzowaną permutację lub właściwy model odpornościowy.
\end{uwaga}

\section{Dane sparowane: test znaków i Wilcoxona}
\begin{twierdzenie}[Test Wilcoxona rangowanych znaków]
\label{thm:wilcoxon-signed-rank}
Niech $D_1,\ldots,D_n$ będą niezależne i mają ten sam
ciągły rozkład symetryczny względem zera, a
$\Prob(D_i=0)=0$. Niech $R_i$ będzie rangą
$|D_i|$ wśród $|D_1|,\ldots,|D_n|$, a
$W_+=\sum_{i=1}^nR_i\ind{\{D_i>0\}}$.
Wówczas
\[
 \E W_+=\frac{n(n+1)}4,\qquad
 \Var(W_+)=\frac{n(n+1)(2n+1)}{24}.
\]
Warunkowo względem wartości bezwzględnych $|D_i|$
rozkład $W_+$ jest taki sam jak rozkład
$\sum_{r=1}^nrB_r$, gdzie $B_r$ są niezależnymi
zmiennymi Bernoulliego o parametrze $1/2$.
\end{twierdzenie}
\begin{proof}
Z ciągłości i symetrii znak każdej różnicy jest
niezależny od jej wartości bezwzględnej, przyjmuje
$+1$ i $-1$ z prawdopodobieństwem $1/2$,
a znaki wszystkich różnic są niezależne.
Po warunkowaniu wartości bezwzględnych rangi tworzą pewną permutację
$1,\ldots,n$, zatem $W_+$ ma wskazany rozkład sumy
niezależnych ważonych zmiennych Bernoulliego.
Jego średnia wynosi $\tfrac12\sum_{r=1}^nr=n(n+1)/4$,
a wariancja $\tfrac14\sum_{r=1}^nr^2
=n(n+1)(2n+1)/24$.
\end{proof}

\begin{uwaga}[Warunki testu i postępowanie z remisami]
Test znaków bada same znaki różnic, podczas gdy
Wilcoxon wykorzystuje dodatkowo rangi wartości
bezwzględnych. Dla różnic o symetrycznym rozkładzie
test Wilcoxona bywa mocniejszy, lecz jego
\emph{dokładny} rozkład opisany wyżej wymaga
braku zerowych różnic i remisów rang.
Przy zerach i remisach trzeba jawnie wybrać
konwencję i kalibrację. Test Wilcoxona nie jest
automatycznie dokładnym testem hipotezy
$\E D_i=0$ dla dowolnego niesymetrycznego rozkładu.
\end{uwaga}

\section{Porównywanie więcej niż dwóch grup i danych sparowanych}
\begin{stwierdzenie}[Statystyka Kruskala--Wallisa i jej kalibracja]
\label{prop:kruskal-wallis}
Niech $k\ge2$ grup składa się z niezależnych próbek,
grupa $j$ ma $n_j\ge1$ obserwacji i
$N=\sum_{j=1}^kn_j$. Przy wspólnym \emph{ciągłym}
rozkładzie pod $H_0$ oznaczmy przez $\overline R_j$
średnią rang w grupie $j$.
Statystyka Kruskala--Wallisa, przy braku remisów, wynosi
\[
 H=\frac{12}{N(N+1)}
 \sum_{j=1}^kn_j
 \left(\overline R_j-\frac{N+1}{2}\right)^2.
\]
Jeżeli $k$ jest stałe, $\min_jn_j\to\infty$ i
$\min_j n_j/N$ pozostaje oddalone od zera,
to $H\Rightarrow\chi^2_{k-1}$ pod $H_0$.
Dla małych prób można zastosować dokładny lub
Monte Carlo test permutacyjny na rangach.
\end{stwierdzenie}
\begin{proof}
Wzór jest definicją statystyki jako ważonej sumy
kwadratów odchyleń średnich rang od rangi globalnej.
Ponieważ pod $H_0$ wszystkie przypisania rang do grup
są jednakowo prawdopodobne, permutacyjna kalibracja
wynika z twierdzenia~\ref{thm:permutation-exact}
(po zastąpieniu dwóch grup dowolnym ustalonym
podziałem na $k$ grup). Asymptotyczny rozkład
$\chi^2_{k-1}$ wynika z wielowymiarowego CLT dla
rang przy losowym podziale i jest tu \emph{przywoływany},
a nie dowodzony w pełnej ogólności.
\end{proof}

\begin{uwaga}[Przykład: trzy grupy i pomiary wielokrotne]
Kruskal--Wallis wymaga niezależnych grup.
Przy pomiarach tego samego pacjenta w kilku warunkach
obserwacje nie są niezależne; stosujemy np.
test Friedmana, który rangowuje warunki \emph{wewnątrz}
każdego bloku lub pacjenta, ewentualnie model mieszany.
W sytuacji z dwoma pomiarami sparowanymi
rozważamy Wilcoxona lub test znaków,
a nie Manna--Whitneya.
Przed wyborem testu zawsze ustalamy, która jednostka
jest niezależna i jaka hipoteza ma zostać sprawdzona.
\end{uwaga}

\section{Test Kołmogorowa--Smirnowa i problem dopasowywania parametrów}
\begin{stwierdzenie}[Bezrozkładowość testu KS dla znanego $F_0$]
\label{prop:ks-distributionfree}
Niech $X_1,\ldots,X_n$ będą niezależne o wspólnej,
\emph{znanej pod hipotezą zerową}, ciągłej i ściśle
rosnącej dystrybuancie $F_0$ na jej nośniku.
Wtedy rozkład statystyki
$D_n=\sup_{t\in\R}|F_n(t)-F_0(t)|$
nie zależy od $F_0$, tylko od $n$.
\end{stwierdzenie}
\begin{proof}
Z transformacji całkowej prawdopodobieństwa
$U_i=F_0(X_i)$ mają niezależny rozkład
$\operatorname{Unif}(0,1)$. Po przepisaniu
$F_n(t)=n^{-1}\sum_{i=1}^n\ind{\{U_i\le F_0(t)\}}$
i przejściu do supremum po $u\in(0,1)$
dostajemy statystykę zbudowaną wyłącznie z
niezależnych zmiennych jednostajnych.
\end{proof}

\begin{uwaga}[Błąd często spotykany przy testach normalności]
Jeśli z tych samych danych oszacujemy średnią i
odchylenie standardowe, a potem obliczymy
$\sup_t|F_n(t)-\Phi((t-\overline X_n)/S_n)|$,
to \emph{nie} wolno używać klasycznych wartości
krytycznych testu KS dla ustalonej, znanej $F_0$.
Rozkład statystyki zmienia się przez estymację parametrów.
W takim problemie rozważamy odpowiednią kalibrację
Lillieforsa, bootstrap parametryczny odtwarzający
\emph{całą} procedurę estymacji albo test Shapiro--Wilka,
którego pełne omówienie umieścimy w osobnym
rozdziale diagnostycznym.
\end{uwaga}

\section{Estymacja gęstości metodą jądrową}
\begin{definicja}[Estymator jądrowy KDE]
\label{def:kde}
Niech $X_1,\ldots,X_n$ będą niezależne o wspólnej
gęstości $f$ względem $\lambda_1$. Niech
$K\colon\R\to[0,\infty)$ będzie mierzalna,
$\int_\R K(u)\,\dd\lambda_1(u)=1$,
a $h>0$ oznacza \emph{szerokość pasma}.
Definiujemy
\[
 \widehat f_{n,h}(x)=\frac1{nh}
 \sum_{i=1}^nK\left(\frac{x-X_i}{h}\right).
\]
\end{definicja}

\begin{twierdzenie}[Obciążenie, wariancja i dobór pasma KDE]
\label{thm:kde-bias-var}
Załóżmy, że $f\in C^2(\R)$, $f$ i $f''$ są ograniczone,
a $f$ jest ciągła w ustalonym $x$.
Niech $K\ge0$ spełnia
$\int_\R K(u)\,\dd\lambda_1(u)=1$, $\int_\R uK(u)\,\dd\lambda_1(u)=0$,
$\mu_2(K)=\int_\R u^2K(u)\,\dd\lambda_1(u)<\infty$,
oraz $R(K)=\int_\R K(u)^2\,\dd\lambda_1(u)<\infty$.
Przy $h=h_n\downarrow0$, $nh_n\to\infty$ mamy
\[
 \begin{aligned}
 \E\widehat f_{n,h}(x)-f(x)
   &=\tfrac12h^2\mu_2(K)f''(x)+o(h^2),\\
 \Var(\widehat f_{n,h}(x))
   &=\frac{f(x)R(K)}{nh}+o((nh)^{-1}).
 \end{aligned}
\]
Jeżeli $f(x)>0$, $f''(x)\ne0$ i $\mu_2(K)>0$,
to minimalizacja głównych członów punktowego MSE daje
$h\asymp n^{-1/5}$, a MSE rzędu $n^{-4/5}$.
\end{twierdzenie}
\begin{proof}
Zmiana zmiennej w całce daje
$\E\widehat f_{n,h}(x)=
\int_\R K(u)f(x-hu)\,\dd\lambda_1(u)$.
Ze wzoru Taylora
$f(x-hu)=f(x)-hu f'(x)+
\tfrac12h^2u^2f''(x)+h^2u^2r_h(u)$,
gdzie $r_h(u)\to0$ dla każdego $u$, a dzięki
globalnej ograniczoności $f''$ mamy jednostajne
ograniczenie $|r_h(u)|\le C$.
Symetria pierwszego momentu jądra i zbieżność
zdominowana dla $u^2K(u)$ dają wzór na obciążenie.
Dla wariancji niezależność implikuje
\[
 \Var(\widehat f_{n,h}(x))=
 \frac1n\left[\frac1h
 \int_\R K(u)^2f(x-hu)\,\dd\lambda_1(u)
 -(\E\widehat f_{n,h}(x))^2\right].
\]
Dominacja przez $\norm f_\infty K(u)^2$ i ciągłość $f$
dają zbieżność pierwszej całki do $f(x)R(K)$.
Drugi składnik jest $O(n^{-1})=o((nh)^{-1})$.
Kwadrat dominującego obciążenia ma rząd $h^4$,
a wariancja rząd $(nh)^{-1}$; zrównanie tych rzędów
daje $h\asymp n^{-1/5}$.
\end{proof}

\begin{uwaga}[KDE przy granicy nośnika]
Ponieważ zwykły estymator KDE nie uwzględnia z góry granic
nośnika, jego poprawność punktowa we wnętrzu przestrzeni nie
musi przenosić się na brzeg.
Jeżeli prawdziwy rozkład ma nośnik $[0,\infty)$,
a stosujemy symetryczne jądro bez korekty, to
przy $x=0$ część masy jądra wychodzi poza nośnik,
a zwykły wynik o obciążeniu rzędu $h^2$ nie musi
obowiązywać. Rozważamy transformację (np. logarytm),
jądro jednostronne lub odpowiednią korektę brzegową.
Przed porównywaniem KDE różnych prób należy też
ustalić sposób doboru pasma: zbyt małe daje poszarpaną
gęstość, zbyt duże ukrywa rzeczywiste wielomodalności.
\end{uwaga}

\section{Estymatory M odporne na obserwacje odstające}
\begin{definicja}[Estymator Hubera w modelu położenia]
\label{def:huber}
Niech $c>0$ i $s>0$ będą ustalonymi parametrami.
Definiujemy funkcję
$\psi_c(u)=\min\{c,\max\{-c,u\}\}$.
\emph{Estymatorem położenia Hubera} nazywamy rozwiązanie
względem $t$ równania
\[
 \sum_{i=1}^n\psi_c\left(\frac{X_i-t}{s}\right)=0.
\]
Jest to warunek pierwszego rzędu minimalizacji
sumy strat $\rho_c((X_i-t)/s)$, gdzie
$\rho_c(u)=u^2/2$ dla $|u|\le c$ oraz
$\rho_c(u)=c|u|-c^2/2$ dla $|u|>c$.
\end{definicja}

\begin{stwierdzenie}[Asymptotyczna wariancja estymatora Hubera]
\label{prop:huber-asympt}
Niech $X_i=\theta+\varepsilon_i$, gdzie $\varepsilon_i$
są niezależne o jednakowym rozkładzie symetrycznym
względem zera, $s>0$ jest znane,
$\E\psi_c(\varepsilon_1/s)^2<\infty$ oraz
$A=\E\psi_c'(\varepsilon_1/s)>0$.
Załóżmy, że rozkład $\varepsilon_1$ nie ma atomów
w punktach $\pm cs$, a rozwiązania $\widehat\theta_n$
są zgodne i spełniają lokalne rozwinięcie
$\sum_i\psi_c((X_i-\widehat\theta_n)/s)
=\sum_i\psi_c(\varepsilon_i/s)
-(\widehat\theta_n-\theta) n A/s
+o_\Prob(\sqrt n+ n|\widehat\theta_n-\theta|)$.
Wówczas
\[
 \sqrt n(\widehat\theta_n-\theta)
 \Rightarrow\Normal\left(0,
 \frac{s^2\E[\psi_c(\varepsilon_1/s)^2]}{A^2}\right).
\]
\end{stwierdzenie}
\begin{proof}
Symetria i nieparzystość $\psi_c$ zapewniają
$\E\psi_c(\varepsilon_1/s)=0$.
CLT daje normalną granicę dla
$n^{-1/2}\sum_i\psi_c(\varepsilon_i/s)$.
Wykorzystujemy rozwinięcie z tezy w równaniu estymatora.
Dodatniość $A$ i zgodność umożliwiają przeniesienie
składnika liniowego na drugą stronę oraz otrzymanie
reprezentacji
\[
 \sqrt n(\widehat\theta_n-\theta)
 =\frac{s}{A\sqrt n}\sum_{i=1}^n
 \psi_c(\varepsilon_i/s)+o_\Prob(1).
\]
Twierdzenie Slutsky'ego prowadzi do wskazanej granicy.
Warunek lokalnego rozwinięcia jest standardowym krokiem
teorii M-estymatorów, a jego weryfikacja dla konkretnego
modelu wymaga odrębnej kontroli regularności.
\end{proof}

\begin{uwaga}[Interpretacja parametru $c$]
Małe $c$ silniej ogranicza wpływ skrajnych reszt,
ale może zmniejszać efektywność w dokładnie normalnym
modelu. Duże $c$ zbliża algorytm do średniej i MNK,
lecz zwiększa wrażliwość na skrajne obserwacje.
Przed zastosowaniem trzeba zdefiniować skalę $s$
(np. odpornie przez MAD), ponieważ próg $cs$ działa
w jednostkach oryginalnych danych. Wyboru $c$
nie należy optymalizować jedynie na podstawie
pojedynczej realizacji z obserwacją odstającą.
\end{uwaga}

\section{Praktyczna mapa decyzji}
\begin{center}\small
\begin{tabular}{>{\raggedright\arraybackslash}p{.26\textwidth}>{\raggedright\arraybackslash}p{.29\textwidth}>{\raggedright\arraybackslash}p{.29\textwidth}}
\toprule
\textbf{Cel} & \textbf{Metoda} & \textbf{Sprawdź najpierw}\\
\midrule
Porównać dwa niezależne rozkłady & Mann--Whitney, KS, permutacje & Czy interesuje Cię położenie, średnia czy cała dystrybuanta?\\
Porównać pomiary przed i po & Wilcoxon, test znaków, model par & Czy parowanie jest rzeczywiste i czy rozkład różnic jest symetryczny?\\
Porównać co najmniej trzy grupy & Kruskal--Wallis & Czy grupy są niezależne, czy to powtarzane pomiary?\\
Porównać powtarzane warunki & Friedman, model mieszany & Pacjent/blok jest jednostką niezależną\\
Zbadać zgodność z zadanym rozkładem & KS, Anderson--Darling, Shapiro--Wilk & Czy parametry hipotezy są znane czy estymowane?\\
Oszacować gęstość i mody & KDE & Pasmo, granice nośnika i liczba obserwacji\\
Oszacować typowy poziom przy outlierach & Mediana, estymator Hubera & Czy celem jest mediana/położenie, czy prawdziwa średnia?\\
\bottomrule
\end{tabular}
\end{center}

\section{Laboratorium: rangi, gęstość i odporność}
\begin{lstlisting}[language=Python,caption={Kompletny program: testy rangowe, KS, KDE oraz estymator Hubera (Python, NumPy i SciPy).}]
import numpy as np
from scipy import stats

x = np.array([9., 10., 10., 11., 12.])
y = np.array([10., 11., 12., 13., 14.])
print("Mann-Whitney:", stats.mannwhitneyu(x, y, alternative="two-sided"))

before = np.array([21., 25., 23., 28., 22., 24.])
after = np.array([20., 23., 20., 24., 17., 18.])
print("Wilcoxon paired:", stats.wilcoxon(after-before, method="exact"))

samples = [np.array([6., 7., 6., 8., 7.]),
           np.array([9., 8., 10., 9., 11.]),
           np.array([11., 12., 10., 13., 12.])]
print("Kruskal-Wallis:", stats.kruskal(*samples))

rng = np.random.default_rng(2026)
z = rng.normal(size=80)
# Valid KS calibration: test against a fully specified N(0,1).
print("KS fixed normal:", stats.kstest(z, "norm", args=(0, 1)))
# For fitted mean and variance, this p-value would not be valid.

contaminated = np.r_[rng.normal(size=80), [15.0, 20.0]]
print("mean, median:", np.mean(contaminated), np.median(contaminated))

scale = stats.median_abs_deviation(contaminated, scale="normal")
c = 1.345
center = np.median(contaminated)
for _ in range(100):
    residual = (contaminated - center) / scale
    weights = np.ones_like(residual)
    mask = np.abs(residual) > c
    weights[mask] = c / np.abs(residual[mask])
    next_center = np.sum(weights * contaminated) / np.sum(weights)
    if abs(next_center - center) < 1e-12:
        break
    center = next_center
print("Huber location:", center)

kde = stats.gaussian_kde(contaminated)
for point in [-1.0, 0.0, 1.0]:
    print("KDE at", point, float(kde(point)[0]))
\end{lstlisting}

\begin{lstlisting}[language=Julia,caption={Kompletny program: odporna estymacja i jądrowa estymacja gęstości (Julia, Random i Statistics).}]
using Random, Statistics
Random.seed!(2026)
x = vcat(randn(80), [15.0, 20.0])
println("mean: ", mean(x), " median: ", median(x))

# Huber location with robust initial MAD scale.
scale = 1.4826 * median(abs.(x .- median(x)))
c = 1.345
center = median(x)
for _ in 1:100
    residual = (x .- center) ./ scale
    weight = min.(1.0, c ./ max.(abs.(residual), 1e-12))
    nextcenter = sum(weight .* x) / sum(weight)
    if abs(nextcenter - center) < 1e-12
        center = nextcenter
        break
    end
    center = nextcenter
end
println("Huber location: ", center)

# Gaussian KDE implemented with elementary functions.
n = length(x)
sigma = std(x)
h = 1.06 * sigma * n^(-1/5)  # rule of thumb; sensitive to outliers
kde(t) = sum(exp.(-0.5 .* ((t .- x)./h).^2)) / (n*h*sqrt(2*pi))
for point in (-1.0, 0.0, 1.0)
    println("KDE at ", point, ": ", kde(point))
end
\end{lstlisting}

\begin{zadanie}[Dlaczego test Manna--Whitneya nie zawsze bada średnią?]
Niech $X$ będzie stała równa zero, a $Y$ przyjmuje
wartość $-1$ z prawdopodobieństwem $0{,}9$ i wartość
$9$ z prawdopodobieństwem $0{,}1$. Wyznaczyć
$\E X$, $\E Y$ oraz $\Prob(X<Y)$.
\par\emph{Rozwiązanie.} Mamy $\E X=0$,
$\E Y=-0{,}9+0{,}9=0$, ale
$\Prob(X<Y)=\Prob(Y=9)=0{,}1$.
Wobec tego równość średnich nie wymusza
$\Prob(X<Y)=1/2$. Test oparty na porządku obserwacji
może wykrywać różnicę mimo równości średnich;
nie jest to sprzeczność, lecz różnica między hipotezami.
\end{zadanie}


